% Generated draft: requires mathematical and source review.
\documentclass[11pt]{article}
\usepackage[margin=25mm]{geometry}
\usepackage{fontspec}
\setmainfont{Latin Modern Roman}
\usepackage{amsmath,amssymb,mathtools,unicode-math}
\setmathfont{Latin Modern Math}
\usepackage{enumitem,needspace,xurl,hyperref}
\setlength{\parindent}{0pt}
\setlength{\parskip}{5pt}
\setlength{\emergencystretch}{3em}
\newcommand{\N}{\mathbb{N}}
\newcommand{\Z}{\mathbb{Z}}
\newcommand{\Q}{\mathbb{Q}}
\newcommand{\R}{\mathbb{R}}
\newcommand{\C}{\mathbb{C}}
\newcommand{\F}{\mathbb{F}}
\newcommand{\A}{\mathbb{A}}
\newcommand{\PP}{\mathbb{P}}
\newcommand{\E}{\mathbb{E}}
\newcommand{\eps}{\varepsilon}
\newcommand{\dd}{\,\mathrm{d}}
\begin{document}
\section*{UnsolvedMath ID: 20003103}
\textbf{Source problem code:} catalogue entry
\subsection{Definitions and mathematical statement}
The reconstructed question is: determine the self quasi-isometry group of Teichm\"uller space equipped with the asymmetric Thurston metric. The appended words ``3. Symmetrization of the Thurston metric'' appear to be the next section heading, not an additional clause of Problem 2.13. The source record is marked reconstructed and unverified, and the complete corrected quotation is absent.

Fix a finite-type surface $S$ in the intended scope, set $d=d_{\mathrm{Th}}$, and put
\[
 D(X,Y)=\max\{d(X,Y),d(Y,X)\}.
\]
This is the max-symmetrization, or length-spectrum metric in the aid's normalization. Define a directional self quasi-isometry to be a map $F$ for which some $K\geq1$, $C\geq0$ satisfy
\[
 K^{-1}d(X,Y)-C\leq d(FX,FY)\leq Kd(X,Y)+C
 \quad\text{for every ordered pair }(X,Y),
\]
and whose image is $C'$-dense for some finite $C'$ in the symmetric metric $D$. Identify two such maps when $\sup_XD(FX,GX)<\infty$. This two-sided equivalence and symmetric coarse-surjectivity give a group under composition. Other asymmetric definitions may not give the same group, so this is an explicit convention rather than a definition supplied by the original sentence.

The aid's reduction is the injective homomorphism
\[
 \operatorname{QI}_{\mathrm{dir}}(\mathcal T(S),d)
       \hookrightarrow\operatorname{QI}(\mathcal T(S),D),
\]
obtained by retaining the same map and equivalence class. Taking maxima proves the symmetric quasi-isometry inequalities, and the chosen bounded-distance equivalence makes injectivity precise.

Its image has a necessary oriented test. For any sequence of ordered pairs $(X_j,Y_j)$ with $D(X_j,Y_j)\to\infty$, a directional quasi-isometry preserves
\[
 \frac{d(X_j,Y_j)}{D(X_j,Y_j)}\longrightarrow0.
\]
That is, forward-sublinear ordered pairs remain forward-sublinear. This is independent of the representative within bounded $D$-distance and does not follow from symmetric distance comparison alone.

The supplied example fixes $X$ and pinches one essential curve $\alpha$ in a family $Y_\eps$, with $\ell_\alpha(Y_\eps)=\eps\to0$, bounded twist, and the remaining Fenchel--Nielsen data controlled in a compact range. It reports
\[
 d(X,Y_\eps)=\log\log(1/\eps)+O(1),\qquad
 d(Y_\eps,X)=\log(1/\eps)+O(1).
\]
Thus $D(X,Y_\eps)=\log(1/\eps)+O(1)$, the thick-to-pinched direction is sublinear relative to $D$, and the reverse direction is dominant. A length-spectrum quasi-isometry lifting to a directional one must preserve this ordered asymptotic distinction. The criterion is necessary, not a classification or a sufficient lifting theorem. The supplied record gives no full determination of either quasi-isometry group.
\subsection{Short English statement}
Determine the quasi-isometry group for the asymmetric Thurston metric, first fixing a group-forming convention. The supplied reduction embeds it in the length-spectrum quasi-isometry group and gives a necessary directional test using one-curve pinching: the two ordered distances grow at different scales. This does not classify the group.
\subsection{Sources}
\begin{itemize}
\item Local UnsolvedMath registry, permanent ID $20003103$, source code catalogue entry; canonical location \texttt{aim-topology-notes.json}, \texttt{notes[190]}.
\item Weixu Su, AIM list \emph{Problems on Thurston Metric}, Problem 2.13, workshop ``Lipschitz metric on Teichmueller space,'' \url{https://aimath.org/pastworkshops/lipschitzteichproblems.pdf}. The extraction appends a symmetrization section heading. The supplied verification reports PDF inspection but the record is marked reconstructed and unverified; no PDF was accessed here.
\item Attached research attempt 1 and summary give the directional reduction and pinching asymptotics without external bibliographic citations. They are marked reduction, proved, and low novelty confidence, with an explicit machine-generated, non-peer-reviewed research-aid caveat requiring independent verification before citation.
\end{itemize}
\end{document}
